\par\textbf{Procedimiento completo para comparar tu solución.} Integrar hidrostática con $T=T_0-\beta h$ y gas ideal da $p=p_0(T/T_0)^{g_0/(R\beta)}$. Después cerrar densidad y calcular sonido para el Mach.

\begin{enumerate}

\item Temperatura local.

\[\fitmath{T=288.15-0.0065\cdot 4500\approx 258.9\,\mathrm{K}}\]

\item Presión de la misma capa.

\[\fitmath{p=101325\cdot {\left(\frac{T}{288.15}\right)}^{\frac{9.80665}{287.05287\cdot 0.0065}}}\]

\[\fitmath{p=101325\cdot {\left(\frac{(258.9)}{288.15}\right)}^{\frac{9.80665}{287.05287\cdot 0.0065}}\approx 57728.3\,\mathrm{Pa}}\]

\item Gas ideal local.

\[\fitmath{\rho=\frac{p}{287.05287\cdot T}}\]

\[\fitmath{\rho=\frac{(57728.30011)}{287.05287\cdot (258.9)}\approx 0.7767743\,\mathrm{kg/m^{3}}}\]

\item Sonido adiabático a temperatura estática.

\[\fitmath{a=\sqrt{1.4\cdot 287.05287\cdot T}}\]

\[\fitmath{a=\sqrt{1.4\cdot 287.05287\cdot (258.9)}\approx 322.5604\,\mathrm{m/s}}\]

\item Mach con la velocidad de ensayo prescrita.

\[\fitmath{\mathrm{Ma}=\frac{140}{a}}\]

\[\fitmath{\mathrm{Ma}=\frac{140}{(322.560356)}\approx 0.4340273\,}\]

\end{enumerate}
